\documentclass[10pt]{article}
\usepackage[margin=0.8in]{geometry}
\usepackage{amsmath,amssymb,hyperref}
\hypersetup{colorlinks=true,urlcolor=blue,linkcolor=blue}
\setlength{\parindent}{0pt}\setlength{\parskip}{4pt}
\begin{document}
\begin{center}
{\Large\bfseries Problem 00000009926: the unrestricted two-invariant clause}\\[5pt]
lizaixi01 \quad AI-led research and formalization \quad 23 October 2026
\end{center}

\textbf{Result and precise limits.} Both source versions assert that every family of two invariants fails to separate some nonisomorphic two-parameter persistence modules, and hence that at least three are needed. No numerical, computable, local, stable, finite-description or codomain restriction is stated. Under this unrestricted assertion, the actual natural-isomorphism-class invariant is already complete. We construct a family of exactly two \emph{different} invariant functions and prove it separates all actual modules, even on the pointwise finite-dimensional subcategory.

The values are isomorphism classes, and the construction is noncomputable. This is not a numerical or algorithmic classification, does not refute impossibility results for discrete complete invariants, and does not establish completeness of the source's rank invariant plus its undefined degeneracy measures. It refutes only the explicit unrestricted necessary lower-bound/nonseparation clause. Additional admissibility restrictions would define another claim and must be supplied before such a lower bound is meaningful.

\section*{1. Genuine persistence modules and natural isomorphisms}
For an actual indexing category $J$, the formal object type is
\[
 \mathcal P_J=\operatorname{Fun}(J,\operatorname{Vect}_{\mathbb Q}).
\]
Lean uses genuine \texttt{Functor}s into \texttt{ModuleCat} over $\mathbb Q$. Morphisms are actual natural transformations and $M\cong N$ is a true categorical natural isomorphism, with linear isomorphisms at every component. The proof is uniform in $J$ and is instantiated with both $(\mathbb N^2,\le)$ and $(\mathbb R^2,\le)$, using their actual coordinatewise preorder categories.

This is the standard functor formulation of multiparameter persistence: Vipond \cite{V}, printed p.~4, defines the coordinatewise indexing categories and vector-space category; Definition~2 on printed p.~5 defines persistence modules as such functors and morphisms as natural transformations. Our field choice $\mathbb Q$ is legitimate for that definition. We use this concrete model directly, without claiming to reprove equivalence with graded modules over polynomial or monoid rings.

The genuine pointwise finite-dimensional condition is
\[
 \operatorname{PFD}(M)\iff
 \forall j\in J,\quad \dim_{\mathbb Q}M(j)<\infty.
\]
Lean encodes it with the actual \texttt{FiniteDimensional} property at each component. No assigned dimension labels or stand-in module objects are introduced.

\section*{2. The actual quotient, with its completeness bridge}
Let $\mathcal Q_J$ be the actual quotient of $\mathcal P_J$ by the equivalence relation
\[
 M\sim N\quad\Longleftrightarrow\quad
 \operatorname{Nonempty}(M\cong N).
\]
The formalization uses Mathlib's \texttt{Skeleton} and its actual natural-isomorphism setoid. Its quotient constructor $q$ satisfies the proved library bridge
\[
 q(M)=q(N)\quad\Longleftrightarrow\quad
 \operatorname{Nonempty}(M\cong N).
\]
This is an actual quotient-equality theorem; the relation is not replaced by equality of a chosen numeric identifier. Consequently $I(M)=q(M)$ is invariant under true natural isomorphisms and complete for them. There is no assertion that $q$ is an effective function on finite presentations.

\section*{3. Two distinct invariants, not repeated indices}
Take the actual constant persistence modules $Z$ and $U$ with components respectively the zero-dimensional rational space $(\operatorname{Fin}0\to\mathbb Q)$ and the one-dimensional space $\mathbb Q$. Both are proved pointwise finite dimensional. For any object $j\in J$, a natural isomorphism $Z\cong U$ would give a genuine linear equivalence
\[
 (\operatorname{Fin}0\to\mathbb Q)\simeq_{\mathbb Q}\mathbb Q.
\]
Its inverse maps $1$ and $0$ into a subsingleton space, so injectivity would force $1=0$ in $\mathbb Q$. Thus $Z$ and $U$ are genuinely nonisomorphic.

Define $K(M)=q(Z)$, a constant invariant. The two functions $I$ and $K$ differ: at $U$, their equality would imply $U\cong Z$, contradicting the component argument. Lean constructs the actual finite set $\mathcal F=\{I,K\}$ and proves
\[
 |\mathcal F|=2.
\]
This is not a two-element indexing type used to repeat the same function. Every member is a genuine isomorphism invariant. If both values agree on $M,N$, in particular $I(M)=I(N)$, so the actual quotient bridge gives $M\cong N$. Conversely a natural isomorphism preserves both values. Hence $\mathcal F$ is a complete family.

\section*{4. Restriction to genuine subclasses and the source negation}
The proof also restricts the same class map to every actual subclass $S\subseteq\mathcal P_J$. For $M,N\in S$, equality of restricted class values is equivalent to a genuine natural isomorphism of their underlying persistence functors. If $S$ contains $U$, this restricted class map still differs from the constant $q(Z)$ map, so the family has cardinality exactly two. In particular, this applies to the actual PFD subclass, since $U$ is proved PFD.

The final theorems for both $\mathbb N^2$ and $\mathbb R^2$ formally negate the statement that every two-element family of isomorphism invariants on PFD persistence modules admits two nonisomorphic modules with all invariant values equal. An explicit complete family of cardinality two contradicts the asserted lower bound three. The actual class map alone is already a complete single invariant; the second distinct function makes the source's exact ``any two'' clause directly testable.

This argument is deliberately about the unrestricted claim as printed. A quotient-valued invariant is tautological from the viewpoint of effective classification, but the source does not exclude it. We do not substitute this object for the rank invariant or claim a usable computation of the proposed degeneracy measures.

\section*{5. Verification and attribution}
All 38 named declarations (21 theorems and 17 definitions/abbreviations) are included in the axiom audit. A new pinned root build, warning-as-error audit and full fresh kernel replay bind the exact source and command logs. Complete original bilingual source, LaTeX, PDF and hash manifests are supplied. No \texttt{sorry}, \texttt{native\_decide} or extra axiom is used. Both PDF pages were rendered and independently checked by worker and supervising AI. No earlier same-ID submission was found in the recorded all-state repository search.

The quotient argument is elementary; the formal work closes the actual persistence, natural-isomorphism, finite-dimensional component, distinct-function and real finite-family cardinality bridges. Research and formalization were led by a Codex worker with separate AI supervision. Organizer acceptance and human verification are not asserted.

\begin{thebibliography}{1}
\bibitem{V} O.~Vipond, \emph{Multiparameter Persistence Landscapes}, Journal of Machine Learning Research 21(61) (2020), 1--38, \href{https://www.jmlr.org/papers/volume21/19-054/19-054.pdf}{official PDF}, printed pp.~4--5, Definition~2. Cited for the actual functor-category model, not for an invariant-classification assertion.
\end{thebibliography}
\end{document}
